\documentclass[11pt]{amsart}

\usepackage[total={6.3in,9.0in},top=1.3in,left=1.1in]{geometry}

\usepackage{enumerate,bm}
\usepackage{palatino}
\usepackage[final]{graphicx}
\usepackage[colorlinks=true,citecolor=blue,linkcolor=red,urlcolor=blue]{hyperref}

\theoremstyle{theorem}
\newtheorem{thm}{Theorem}
\newtheorem{Cor}{Corollary}
\newtheorem{lem}{Lemma}
\newtheorem{lem*}{Lemma}

\theoremstyle{definition}
\newtheorem{def*}{Definition}
\newtheorem{ex*}{Example}

\newcommand{\Div}{\nabla\cdot}
\newcommand{\grad}{\nabla}

\newcommand{\RR}{\mathbb{R}}

\newcommand{\bn}{\mathbf{n}}
\newcommand{\bu}{\mathbf{u}}
\newcommand{\bv}{\mathbf{v}}
\newcommand{\bx}{\mathbf{x}}
\newcommand{\by}{\mathbf{y}}
\newcommand{\bz}{\mathbf{z}}

\newcommand{\bpi}{\bm{\pi}}
\newcommand{\bzero}{\bm{0}}


\title[Reynold's transport theorem for glaciers]{Reynold's transport theorem for glaciers, \\ with applications}

\author{Ed Bueler}

\author{Romain Hugonnet}

\begin{document}

\begin{abstract}
We consider the form that Reynold's transport theorem \cite{Holmes2019} takes for time-dependent 3D sets defined by a fixed lower surface $z=b(x,y)$ and a time-dependent upper surface $z=s(t,x,y)$.  The map-plane points $(x,y)$ cover a fixed domain $\Omega\subset \RR^2$, a partially-glaciated region, and ice exists where $b(x,y)<z<s(t,x,y)$.  The 3D sets considered here, namely time-dependent domains of glacier ice, can have any outline or connectedness in their map-plane projection; this projection is the glacier extent within $\Omega$.  We assume that the 3D boundary of the glacier ice domain consists, up to a set of measure zero, only of a lower and an upper surface; we neglect any cliffs along the margins of the glaciers.  The transport theorem is primarily applied to the evolution of total glacier mass and mass changes for general mass density fields.  We also consider enthalpy, age, and linear momentum as examples.
\end{abstract}

\maketitle

\section{Basic theory} \label{sec:basic}

\subsection{Geometrical formulas for an evolving regional glacier domain}  Consider a fixed map-plane domain $\Omega \subset \RR^2$, an open and connected set.  This set is often something simple like a rectangle around a glaciated region.  All calculations to follow are defined within this set, or on subsets of it, or on three-dimensional sets defined over it.

Suppose that the bed elevation of the glacier(s) is given by a continuous, time-independent function $b(x,y)$ defined on $\Omega$.  All formulas for the general case $b(t,x,y)$ are given at the end, in Appendix \ref{app:general}.

The upper surface elevation $s(t,x,y)$ is free to change in time.  It is defined for any $(x,y)\in \Omega$, and for every $t$ in some duration $[0,T]$ of interest, and it is assumed to be continuous.  We will use $s$ in a manner which is compatible with digital elevation models.  That is, there is a single surface elevation on $\Omega$, which is the glacier surface on ice-covered land, but which is otherwise equal to the bed elevation.  That is, on ice-free land we have $s(t,x,y) = b(x,y)$.

Now define the time-dependent open, three-dimensional subset which is the ice domain:
\begin{equation}
\Lambda(t) = \{(x,y,z) : b(x,y) < z < s(t,x,y)\} \subset \RR^3  \label{eq:Lambda}
\end{equation}
Furthermore define the two-dimensional set
\begin{equation}
\Pi(t) = \{(x,y) : b(x,y) < s(t,x,y) \} \subset \Omega.  \label{eq:Pi}
\end{equation}
This is the open time-dependent horizontal footprint of the glacier, the ice-covered map-plane domain on which the glacier has positive thickness $h = s-b > 0$.  We might write $\Pi(t) = \Lambda(t)^\top$, the (open) vertical projection of $\Lambda(t)$, because $\Pi(t)$ is the set of $(x,y)$ such that there exists $z$ satisfying $(x,y,z) \in \Lambda(t)$.  Note that $\Pi(t)$ is not necessarily connected or simply-connected.

\begin{figure}
FIXME: SHOW CARTOON OF $\Omega$, $\Lambda(t)$, $\Pi(t)$, $\gamma(t)$
\label{fig:setscartoon}
\caption{Over a fixed map-plane region $\Omega$ we define three time-dependent sets for the glaciers: the 3D ice domain $\Lambda(t)$, its map-plane projection $\Pi(t)$, and the map-plane curve $\gamma(t)$ for the margin of the glaciers.}
\end{figure}

Finally, let
\begin{equation}
\gamma(t) = \partial \Pi(t) \subset \bar\Omega.  \label{eq:gamma}
\end{equation}
This is the map-plane projection of the marginal curve surrounding the glacier(s) at time $t$.  Along $\gamma(t)$ the thickness vanishes because, by continuity, $s=b$ at points of $\gamma(t)$.  We will use the fact that $\gamma(t)$ is a measurable set, which follows from continuity of $b$ and $s$.  One would suppose that $\gamma(t)$ is one-dimensional, and thus has measure zero in $\Omega$ at every time, so that $\int_{\gamma(t)} f\,dx\,dy = 0$ for any function $f$, but we will not actually need this.

We may write the surface using a level set function,
	$$f(t,x,y,z) = z - s(t,x,y),$$
so that the upper surface of the glacier(s) is $f=0$, that is, $z=s(t,x,y)$.  From now on we will assume that $f$ and its partial derivatives exist and are integrable; equivalently, we make these assumption about $s$.)  Then $\nabla f = (-s_x,\,-s_y,\,1)$ and $|\nabla f| = \sqrt{1+s_x^2+s_y^2}$.

Because the gradient of $f$ is orthogonal to its level surfaces, the upward unit normal vector on the upper surface of the glacier is $\hat \bn_s = \grad f/|\grad f|$ or
\begin{equation}
\hat \bn_s = \frac{(-s_x,\,-s_y,\,1)}{\sqrt{1+s_x^2+s_y^2}}.  \label{eq:hns}
\end{equation}
The un-normalized vector will also be used:
\begin{equation}
\bn_s = \grad f = (-s_x,\,-s_y,\,1).  \label{eq:ns}
\end{equation}

Consider the (scalar) velocity of the surface in the upward norm direction.  To derive a formula, let $\bx(t)$ be a differentiable path tracking the surface, i.e.~of a non-material (notional) point that is always on the surface.  It satisfies $f(t,\bx(t)) = 0$ and thus $f_t + \nabla f \cdot \bx' = 0$ by the chain rule.  We may split $\bx' = v_n \hat \bn_s + \bv_\parallel$ with $v_n$ a scalar and with the vector $\bv_\parallel$ tangential to the surface.  However, since $\nabla f$ is normal to the surface, $\nabla f\cdot \bv_\parallel = 0$, so from $\hat \bn_s = \grad f/|\grad f|$ we have $\grad f \cdot \hat \bn_s = |\grad f|$.  This leaves $f_t + |\nabla f|\,v_n = 0$.  It follows that $v_n = -f_t/|\nabla f|$, and thus the desired normal velocity of the surface is
    $$v_n = \frac{s_t}{\sqrt{1+s_x^2+s_y^2}}.$$
Observe that the (vector) tangential portion of the surface velocity, namely $\bv_\parallel$, involves an arbitrary parameterization choice, but we will not need it anyway.

\subsection{Reynold's transport theorem for regional glacier geometry}  Suppose $\phi(t,x,y,z)$ is the density of some physical quantity, in quantity-per-volume units, which is being transported by the ice flow.  The density could be a scalar, vector, or tensor field.  For example, it might be the (variable) ice mass density $\rho$, in mass-per-volume units (Section \ref{sec:massapp}).  On the other hand, perhaps $\phi = \rho \psi$ for some specific quantity $\psi(t,x,y,z)$ in per-mass units.  For example $\psi$ could be (specific) enthalpy \cite{Aschwandenetal2012} in energy-per-mass units so that $\phi$ is an energy density with energy-per-volume units (Section \ref{sec:otherapp}).

Define
\begin{equation}
F(t) = \int_{\Lambda(t)} \phi(t,x,y,z) \, dx\,dy\,dz, \label{eq:F}
\end{equation}
or $F(t)=\int_{\Lambda(t)} \phi \, dV$ for short, where $dV=dx\,dy\,dz$.  Thus $F(t)$ is the total amount of $\phi$ which is within the time-dependent ice domain $\Lambda(t)$.  Our goal is to compute $F'(t)$, and understand the various contributions to it.

We may apply Fubini's theorem to the three-dimensional integral which defines $F$.  Let
\begin{equation}
g(t,x,y) = \int_{b(x,y)}^{s(t,x,y)} \phi(t,x,y,z) \, dz, \label{eq:g}
\end{equation}
so that
	$$F(t) = \int_{\Pi(t)} g(t,x,y) dx\,dy,$$
or $F(t)=\int_{\Pi(t)} g \, dA$ for short, where $dA=dx\,dy$.  (For comparison, $A(t) = \int_{\Pi(t)} 1\,dA$ is the glaciated area within $\Omega$.)  Note that $g$ can only be nonzero at $x,y \in \Pi(t)$, that is, it can only be nonzero in ice-filled columns where $s$ exceeds $b$, so we declare that $g(t,x,y)=0$ for $(x,y) \notin \Pi(t)$, which is compatible with its integral definition \eqref{eq:g}.  In fact, let us make explicit our hypothesis of admissibility, namely $s\ge b$, as well as our assumption of continuity for $s$ and $b$.  Thus $g$ is defined for all $(x,y)\in\Omega$.

To find the ordinary derivative of $F$ we will need to differentiate $g$ with respect to $t$ (partial derivative), which we do using the familiar Leibniz rule.  Since $b_t=0$ (versus in Appendix \ref{app:general}), only the upper limit contributes a boundary term:
    $$g_t(t,x,y) = \int_{b(x,y)}^{s(t,x,y)} \phi_t(t,x,y,z)\,dz + \phi(t,x,y,s(t,x,y))\, s_t(t,x,y).$$
This formula applies everywhere on $\Omega$, but note that at points $(x,y)$ outside of $\Pi(t)$ we have $\int_b^s f \,dz=0$ (any $f$) and $s_t=0$.  Evidently, $g_t$ is zero outside of $\bar\Pi(t)$.

Now the main idea is that the map-plane ice-covered domain $\Pi(t) \subset \Omega$ changes with time, and it has a moving boundary $\gamma(t)$.  The 2D form of Reynold's transport theorem \cite{Holmes2019} applies to $g$ over $\Pi(t)\subset \Omega$, yielding
\begin{equation}
\frac{d}{dt}\left(\int_{\Pi(t)} g(t,x,y)\,dx\,dy\right) = \int_{\Pi(t)} g_t(t,x,y)\,dx\,dy \;+\; \underbrace{\int_{\gamma(t)} g\, u_\gamma \, d\ell}_{\ast},  \label{eq:reynoldstwod}
\end{equation}
where $u_\gamma$ denotes the outward normal velocity of the margin curve $\gamma(t)$ within the $(x,y)$-plane, i.e.~within $\Omega$.  It is a scalar function along this curve.\footnote{Again we observe that the scalar velocity $u_\gamma$---we might say ``speed'' but it can be of either sign---is well-defined even though the full vector velocity of $\gamma$ is not meaningful.  The latter would require a parameterization choice.}

The line integral $\ast$ over $\gamma(t)$ is important to consider, but in fact it vanishes because the value of $g$ along $\gamma(t)$ is zero.  That is, $g(t,x,y)\big|_{\gamma(t)} = 0$ because $s=b$ along $\gamma(t)$.  The measure of $\gamma(t)$ is thus not important.  The reason the line integral vanishes is the same mechanism as the following 1D fact: if $F(t)=\int_{a(t)}^{b(t)} g(t,x)\,dx$, with $g$ vanishing at both endpoints for all $t$, namely $g(t,a(t))=g(t,b(t))=0$, then $F'(t) = \int_{a(t)}^{b(t)} g_t(t,x)\,dx$ from the Leibniz rule, with no boundary-motion terms, even though $a(t),b(t)$ move.  The vanishing of the integrand at the moving boundary kills the boundary-motion contribution.

For the above logic to be valid, $\phi(t,x,y,z)$ must stay bounded as we approach $\gamma(t)$.  Formally, for fixed $t_0$, if we parameterize the curve $\gamma(t_0)$ by $(x(\tau), y(\tau))$ and choose any parameter value $\tau_0$ then $\phi$ must be well-behaved enough so that
    $$\lim_{(x,y) \to (x(\tau_0), y(\tau_0))} g(t_0,x,y) = \lim_{(x,y) \to (x(\tau_0), y(\tau_0))} \int_{b(x,y)}^{s(t_0,x,y)} \phi(t_0,x,y,z) \, dz = 0.$$
Otherwise the above argument which justifies the vanishing of the line integral would be defeated by an integrand which blows up competitively.  An integrability assumption on $\pi$ is needed; an assumption of boundedness or continuity certainly suffices but it need not be that strong.

At this point we have derived a useful result, namely a form of Reynold's transport theorem \cite{Holmes2019} suited to continuum discussions of glaciers.

\begin{thm}  \label{thm:newreynolds}  Suppose the quantity $\phi(t,x,y,z)$ is defined at points of $\Lambda(t)$, and integrable, and regular enough to have a surface trace $\phi|_s(t,x,y) = \phi(t,x,y,s(t,x,y))$ along $z=s(t,x,y)$.  Then
\begin{align*}
\frac{d}{dt}\left(\int_{\Lambda(t)} \phi(t,x,y,z)\,dx\,dy\,dz\right) &= \int_{\Lambda(t)} \phi_t(t,x,y,z)\,dx\,dy\,dz \\
   &\qquad + \int_{\Pi(t)} \phi|_s(t,x,y)\, s_t(t,x,y)\,dx\,dy.
\end{align*}
With simpler notation, using definition \eqref{eq:F} and suppressing dependencies, we have
\begin{equation}
F'(t) = \int_{\Lambda(t)} \phi_t\,dV + \int_{\Pi(t)} \phi|_s \, s_t\,dA, \label{eq:newreynolds}
\end{equation}
where $dV=dx\,dy\,dz$ and $dA=dx\,dy$.
\end{thm}

\begin{ex*}  As a simplest example of Theorem \ref{thm:newreynolds}, suppose $\phi=1$ identically.  Noting that
\begin{equation}
V(t) = \int_{\Lambda(t)} 1\,dV \label{eq:V}
\end{equation}
is the total ice volume, and that the first integral vanishes ($1_t=0$), we find
\begin{equation}
V'(t) = \int_{\Pi(t)} s_t\,dA. \label{eq:dVdt}
\end{equation}
This holds because the bed elevation is time-independent, and, critically, because no volume flux is allowed to sneak through at the glacier margins.
\end{ex*}

Despite the fact that $\Pi(t)$ changes in time, and even though the glacier margins are moving, our form of Reynold's transport theorem is actually the same as what holds for a glacier constrained by walls and subject to $s>b$ everywhere in $\Pi(t)$, as used in reference \cite{LofgrenAhlkronaHelanow2022} for example.  In Theorem \ref{thm:newreynolds} the footprint $\Pi(t)$ is treated simply, as if it were not moving.  The Theorem is also compatible with vertically-integrated conservation laws, e.g.~equations of thickness evolution for a layer \cite{Bueler2021conservation}, but that alternative parameterization of free-margin glacier geometry makes an assumption of incompressibility or constant density.

Note that the assumption of time-independent bed elevation ($b_t=0$) is important in the statement of the theorem, but it is generalized in Appendix \ref{app:general} to allow a time-dependent bed elevation.  However, the statement of Theorem \ref{thm:newreynolds} is correct as stated even if there is basal sliding or a basal melt rate.  All we have assumed is that the lower surface of $\Lambda(t)$ \emph{is not moving}; this does not represent an assumption on the basal value of the ice density or of the ice velocity.


\section{Mass applications}  \label{sec:massapp}

The above form of Reynold's transport theorem relates the rate of change of the total of a quantity, within a fixed map-plane region $\Omega$, to observable surface fields even as the thickness \emph{and extent} of the glaciers change.  Thus, of the two integrals in \eqref{eq:newreynolds}, the one which is over $\Pi(t) \subset \Omega$ is often of greatest interest in the mass applications below.  In this section we apply Theorem \ref{thm:newreynolds} to mass transport within a glaciated region, for a completely-general ice mass density field.

\subsection{Total glacier mass, for variable ice density}  \label{ssec:mass} Suppose $\phi(t,x,y,z)=\rho(t,x,y,z)$ is the mass density of the glacier(s), in SI units $\text{kg}\,\text{m}^{-3}$.  We assume this field is measurable and bounded; there is a maximum density $\rho_0$ so that $0 < \rho(t,x,y,z) \le \rho_0$ at all points $(x,y,z) \in \Lambda(t)$.

The total mass of the glaciers in $\Omega$ is
\begin{equation}
M(t) = \int_{\Lambda(t)} \rho(t,x,y,z)\,dx\,dy\,dz. \label{eq:M}
\end{equation}
Applying Theorem \ref{thm:newreynolds} gives
\begin{equation}
M'(t) = \int_{\Lambda(t)} \rho_t\,dV + \int_{\Pi(t)} \rho|_s\,s_t\,dA, \label{eq:dMdt}
\end{equation}
where $\rho|_s(t,x,y) = \rho(t,x,y,s(t,x,y))$ is the surface trace of $\rho$.

Note that no mass flux is permitted at the glacier margins.  Also, many glacier modeling treatments would lack the first integral because of an assumption of constant ice density, not made here, but in fact we see below that this $\rho_t$ term can be absorbed into boundary terms.

Each integral in \eqref{eq:dMdt} could be negative, zero, or positive.  There are many possibilities, but as illustration let us give two cases where the integrals come out zero.  The first integral could be zero if some part of the glacier was densifying ($\rho_t>0$) and an equal amount was becoming less dense ($\rho_t<0$).  The second integral could be zero if the density at the surface was constant but the glacier surface was raising ($s_t>0$) and lowering ($s_t<0$) in different map-plane areas.

\subsection{Free-surface and mass-continuity equations}  \label{ssec:skemc}  To further study mass transport we recall equations for how the mass density evolves and the glacier surface moves.  Both equations follow from the principle of mass conservation.

First, the variable density field $\rho(t,x,y,z)$ satisfies the mass-continuity equation \cite{GreveBlatter2009}:
\begin{equation}
\rho_t + \Div\left(\rho\bu\right) = 0. \label{eq:mc}
\end{equation}
Here $\bu(t,x,y,z)$ is the ice flow velocity, defined within $\Lambda(t)$.  This equation applies everywhere in the 3D set $\Lambda(t)$.  Inserting this into \eqref{eq:dMdt} and applying the divergence theorem we get
\begin{equation}
M'(t) = - \int_{\partial\Lambda(t)} \rho \bu\cdot \hat\bn\,dS + \int_{\Pi(t)} \rho|_s\,s_t\,dA. \label{eq:dMdt:mc}
\end{equation}
Here $\hat\bn$ is the unit normal vector on the boundary surface $\partial\Lambda(t)$, and $dS$ denotes its surface measure.  Equation \eqref{eq:dMdt:mc} will be updated below after we consider the upper surface.

Suppose we are given a time- and space-dependent surface mass balance (SMB) rate function $a(t,x,y)$, provided in mass flux units, namely $\text{kg}\,\text{m}^{-2}\,\text{s}^{-1}$ in SI.  This is the climate experienced by our land-based glaciated region.  By definition, the SMB is the balance between accumulating snow and the loss of melt water, through runoff, measured vertically at the upper surface of the glacier \cite{Cogleyetal2011}.

Because a simulated glacier needs to be able to advance into unglaciated locations, in a time-dependent (prognostic) model we would assume that the SMB $a$ is defined everywhere in $\Omega$, whether ice-covered or not.  In ice-free areas $a$ should have the value which an ice surface at elevation $b$ would experience.  This can be computed from precipitation and an energy balance model \cite{GreveBlatter2009} by snow accumulation minus the ablation modeled using the available energy for melt (if ice were present).  For ice-free areas we expect $a<0$; if $a$ is continuous and it ever becomes positive at some location in $\Omega$ then there is a glacier there \cite[for example]{JouvetBueler2012}.

The motion of the glacier surface is determined by a combination of the flow and the SMB.  To state this, denote the surface trace of the mass density by $\rho|_s(t,x,y) = \rho(t,x,y,s(t,x,y))$, and of the ice velocity by $\bu|_s(t,x,y) = \bu(t,x,y,s(t,x,y))$.  Using definition \eqref{eq:ns}, the free-surface equation is
\begin{equation}
\rho|_s \left(s_t - \bu|_s \cdot \bn_s\right) = a,  \label{eq:ske}
\end{equation}  
which says that the (non-material) ice surface moves vertically according an ice velocity surface component and the SMB.  It is a standard description of glacier geometry evolution \cite{GreveBlatter2009,SchoofHewitt2013}, valid regardless of the density.  It is equivalent to the more familiar vertically-integrated mass-continuity equation for thickness \cite{Bueler2021conservation,GreveBlatter2009,JouvetBueler2012} under an additional incompressibility (constant density) assumption which we do not make here.  Though equation \eqref{eq:ske} is often described as merely kinematical, it is really an important statement of mass conservation \cite[Section 3.3]{Aschwandenetal2012}.

Returning to the rate of change $M'(t)$ of the total glacier mass, free-surface equation \eqref{eq:ske} allows us to rewrite the second integral in \eqref{eq:dMdt} and \eqref{eq:dMdt:mc} as two terms, and one of them is true surface flux.  To see this, define the curved upper surface of the glacier, as a set in 3D, by
\begin{equation}
\Gamma(t) = \{(x,y,z)\,:\,z = s(t,x,y)\} \subset \partial\Lambda(t),  \label{eq:Gamma}
\end{equation}
and denote the element of surface area on $\Gamma(t)$ by $dS = \sqrt{1+s_x^2+s_y^2}\;dx\,dy$.  Recalling definitions \eqref{eq:hns} and \eqref{eq:ns}, equation \eqref{eq:ske} yields
\begin{equation}
\int_{\Pi(t)} \rho|_s\,s_t\,dA = \int_{\Pi(t)} \left(a + \rho|_s\,\bu|_s \cdot \bn_s\right)\,dA = \int_{\Pi(t)} a \,dA + \int_{\Gamma(t)} (\rho\bu)|_s \cdot \hat \bn_s\,dS
\label{eq:Gammamassflux}
\end{equation}
Writing the integral of $a$ over $\Pi(t)$ makes sense because $a$ is measured vertically, whereas the final integral is the classical ice-velocity mass flux through the (non-material) upper surface $\Gamma(t)$, a surface living in 3D space.

Combining the last form in \eqref{eq:Gammamassflux} with equation \eqref{eq:dMdt:mc}, the upper surface flux cancels,\footnote{This cancellation would be made more complicated by allowing mass flux through the glacier margins, which we do not do in this work.} leaving only the mass flux on the base.  Let
\begin{equation}
\Xi(t) = \{(x,y,z)\,:\,z = b(x,y)\} \subset \partial\Lambda(t), \label{eq:Xi}
\end{equation}
be the lower surface of the ice.  Note that it is indeed a time-dependent set, but only because the glacier-covered area can change.  (The vertical projection of $\Xi(t)$ into the map-plane is $\Pi(t)$, which changes in time.)  Then
\begin{equation}
M'(t) = \int_{\Pi(t)} a \,dA - \int_{\Xi(t)} (\rho\bu)|_b \cdot \hat \bn_b\,dS \label{eq:dMdt:best}
\end{equation}
were $dS = \sqrt{1+b_x^2+b_y^2}\;dx\,dy$ and $\hat \bn_b = (b_x, b_y, -1) / \sqrt{1+b_x^2+b_y^2}$ is the outward, i.e.~downward, unit normal on the underside of the 3D ice domain $\Lambda(t)$.

The last integral in \eqref{eq:dMdt:best} is an advective mass flux, caused by the ice flow velocity field $\bu$, through the glacier's curved non-material lower surface $\Xi(t)$ in 3D.  Physically, because the (non-material) basal surface is not moving, this corresponds to basal melt or refreeze.  If desired the integral can be written over the map-plane region instead: $\int_{\Pi(t)} (\rho\bu)|_b \cdot \bn_b\,dA$.

Neither equation \eqref{eq:dMdt} or \eqref{eq:dMdt:best} is new, but they give different views of regional mass (im)balance, the rate of change in the total ice mass in the glaciated region, based on clear mechanisms.  For example, \eqref{eq:dMdt:best} says that total mass changes only by\renewcommand{\labelenumi}{\roman{enumi}.}
\begin{enumerate}
\item $\int_{\Pi(t)} a \,dA$, the total climatic influence of the SMB, minus
\item $\int_{\Xi(t)} (\rho\bu)|_b \cdot \hat \bn_b\,dS$, the total outward (downward) mass flux through the curved lower surface of the glacier.
\end{enumerate}
However, even in the satellite era, the basal value of the ice velocity ($\bu|_b$) of a land-based glacier can hardly be regarded as a standard observable.

\subsection{Effective density of a mass change}  \label{ssec:effdensity} FIXME Denote the overall (average) glacier density at time $t$ by\footnote{This lumped (i.e.~whole glacier) quantity is denoted ``$\rho$'' in \cite{Huss2013} and \cite{HussHugonnet2026}.  Likewise, our quantity $R_\Delta$ is denoted ``$\rho_{\Delta V}$'' there.  However, here we reserve the symbol $\rho$ for the extensive/local density $\rho(t,x,y,z)$ already introduced.}
\begin{equation}
R(t) = \frac{M(t)}{V(t)}. \label{eq:R}
\end{equation}
This is well-defined for any glacier or glaciated region whose total volume function $V(t)$ remains bounded below by a positive constant.

Next, following \cite{HussHugonnet2026}, define the effective density of a mass change during the time interval between $t$ and $t+\Delta t$ as $R_{\Delta} = \Delta M/\Delta V$ or, more precisely,
\begin{equation}
R_{\Delta}(t;\Delta t) = \frac{M(t+\Delta t) - M(t)}{V(t+\Delta t) - V(t)}.  \label{eq:RDeltaearly}
\end{equation}
This effective density has a small-interval limit as long as $M(t)$ and $V(t)$ are differentiable.  Using the mean value theorem we have
\begin{align*}
\lim_{\Delta t\to 0} R_{\Delta}(t;\Delta t) &= \lim_{\Delta t\to 0}  \frac{M(t+\Delta t) - M(t)}{V(t+\Delta t) - V(t)} = \lim_{\Delta t\to 0} \frac{M'(t)\Delta t + o(\Delta t)}{V'(t)\Delta t + o(\Delta t)} = \frac{M'(t)}{V'(t)}.
\end{align*}
(Here $o(\Delta t)$ denotes a quantity with $\lim_{\Delta t \to 0} o(\Delta t)/\Delta t = 0$.)  With slight abuse of notation, we therefore define
\begin{equation}
R_{\Delta}(t) = \frac{M'(t)}{V'(t)}  \label{eq:RDelta}
\end{equation}
so that $R_{\Delta}(t;\Delta t) = R_{\Delta}(t) + o(\Delta t)$.

According to \cite{Huss2013,HussHugonnet2026}, glacier mass changes at global scale can be best informed by observations using estimates of such effective densities, which are evaluated regionally.  Our goal for now is to express the effective density $R_\Delta$ using Reynold's theorem, to clarify its meaning.

First, reference \cite[equation (4)]{Huss2013} states
\begin{equation}
R_{\Delta} = V \frac{\Delta R}{\Delta V} + R. \hspace{7mm} (\emph{\text{approximate}})  \label{eq:RDeltaHuss}
\end{equation}
This turns out to be true up to lower-order terms $o(\Delta t)$ using our definitions above; see Appendix \ref{app:hussverif}.  This form can be checked using the mean value theorem or equivalent; only the ordinary derivatives of $M(t)$ and $V(t)$ are involved in the verification.

By Theorem \ref{thm:newreynolds} the effective density $R_\Delta$ can be computed in terms of spatially-distributed density and geometry changes for a well-observed glaciated region.  In particular, starting from definition \eqref{eq:RDelta}, equations \eqref{eq:dVdt}, \eqref{eq:dMdt}, and \eqref{eq:dMdt:ske} yield these forms for the effective density:
\begin{align}
R_\Delta(t) &= \frac{\int_{\Lambda(t)} \rho_t\,dV + \int_{\Pi(t)} \rho|_s\,s_t\,dA}{\int_{\Pi(t)} s_t\,dA} \label{eq:RDelta:reynolds} \\
    &= \frac{\int_{\Lambda(t)} \rho_t\,dV + \int_{\Pi(t)} a \,dA + \int_{\Gamma(t)} (\rho\bu)|_s \cdot \hat \bn_s\,dS.}{\int_{\Pi(t)} s_t\,dA} \notag
\end{align}

Note that there is no apparent conceptual benefit of rewriting the denominators in \eqref{eq:RDelta:reynolds} using the free-surface equation; we may naturally regard $V'(t)$ as purely kinematical.

Of the terms in the numerators in \eqref{eq:RDelta:reynolds}, presumably a firn densification model would inform the first integral over $\Lambda(t)$, and perhaps only over that portion of $\Lambda(t)$ which is close to the surface.  The integral of the SMB $a$ over $\Pi(t)$ would be modeled as usual by precipitation and melt processes.  The last integral in the numerator, over $\Gamma(t)$, uses the 3D velocity of the ice at the ice surface, presumably modeled by a stress balance for the glacier ice.  However, the vertical component of the surface velocity, i.e.~$\bu|_s\cdot \hat\bz$, is generally hard to observe and/or separate from the geometrical rate $s_t$.


\section{Other applications}  \label{sec:otherapp}

\subsection{Enthalpy.}  \label{ssec:enthalpy}  Suppose the ice is regarded as a mixture of solid-phase ice and liquid water.  Let $\phi=\rho \eta$ where $\eta$ is the (specific) internal energy of the mixture.  Because the phases have different densities and the mixture fraction varies in space and time, the ice is now a variable-density fluid; it makes up temperate ice and polythermal glaciers \cite{GreveBlatter2009}.  The mixture theory which supports these concepts is laid out for glaciers in \cite{Aschwandenetal2012}.

The internal energy $\eta$, from now on called the (specific) enthalpy, is the energy excluding only the kinetic energy.  It has units of energy-per-mass, $\text{J}\,\text{kg}^{-1}$ in SI, so $\phi=\rho\eta$ has units of energy-per-volume ($\text{J}\,\text{m}^{-3}$).

Consider the total enthalpy (internal energy) $E(t)=\int_{\Lambda(t)} \rho\eta\,dV$.  Applying Theorem \ref{thm:newreynolds} and free-surface equation \eqref{eq:ske} we find
\begin{equation}
E'(t) = \int_{\Lambda(t)} (\rho\eta)_t\,dV + \int_{\Pi(t)} a\eta|_s \,dA + \int_{\Gamma(t)} (\rho\eta\bu)|_s \cdot \hat \bn_s\,dS.  \label{eq:dEdt}
\end{equation}

FIXME talk

\subsection{Ice age.}  \label{ssec:age}  FIXME

\subsection{Linear momentum.}  \label{ssec:momentum}  Finally, suppose that the density $\phi$ is the vector-valued linear momentum, namely $\bpi = \rho \bu$.  Then  FIXME


\section{Summary: regional views of total quantities}

The main advantage of using Reynold's transport theorem in the new form, Theorem \ref{thm:newreynolds}, is that the rate of change of the total of a quantity, defined across a fixed map-plane region, can be computed by integrating well-defined spatially-distributed fields even as the glaciated area changes in time.


\bibliography{layer}
\bibliographystyle{siamplain}

\appendix

\section{Verification of a formula from \cite{Huss2013}}  \label{app:hussverif}

Recall definitions \eqref{eq:R} and \eqref{eq:RDeltaearly}.  This Appendix gives a brief calculation showing why our equation \eqref{eq:RDeltaHuss}, which is stated as exact in \cite{Huss2013}, is at least correct up to lower order terms.  That is, we verify that
\begin{equation}
V(t)\, \frac{R(t+\Delta t) - R(t)}{V(t+\Delta t) - V(t)} + R(t) = R_{\Delta}(t) + o(\Delta t).  \label{eq:fromHuss}
\end{equation}
We assume that $V(t)$ and $M(t)$ are continuous and differentiable.

\begin{proof}  Find common denominators, and expand $t+\Delta t$ values using the mean value theorem, or equivalently Taylor's theorem with remainder.  Start from the fraction on the left of \eqref{eq:fromHuss}:
\begin{align*}
V(t)\, \frac{R(t+\Delta t) - R(t)}{V(t+\Delta t) - V(t)} &= V(t)\, \frac{\frac{M(t+\Delta t)}{V(t+\Delta t)} - \frac{M(t)}{V(t)}}{V'(t) \Delta t + o(\Delta t^2)} \\
    &= \frac{V(t) M(t+\Delta t) - V(t+\Delta t) M(t)}{V'(t) \Delta t \, V(t+\Delta t) + o(\Delta t^2)} \\
    &= \frac{V(t) \left[M(t+\Delta t)-M(t)\right] - \left[V(t+\Delta t) - V(t)\right] M(t)}{V'(t) \Delta t \, V(t) + o(\Delta t^2)} \\
    &= \frac{V(t) M'(t)\Delta t - V'(t)\Delta t M(t) + o(\Delta t^2)}{V'(t) \Delta t \, V(t) + o(\Delta t^2)} \\
    &= \frac{V(t) M'(t)}{V'(t) \, V(t) + o(\Delta t)} - \frac{V'(t) M(t)}{V'(t) \, V(t) + o(\Delta t)}  + o(\Delta t)\\
    &= \frac{M'(t)}{V'(t)} - \frac{M(t)}{V(t)}  + o(\Delta t) = R_{\Delta}(t) - R(t) + o(\Delta t). \qedhere
\end{align*}
\end{proof}


\section{Formulas with time-dependent bed elevation}  \label{app:general}

FIXME

\end{document}
